\section{Global coefficient filters and a length-valued Euler product}
\label{sec:globalfilters}

Continue throughout with a closed hyperbolic surface of genus at least
two; the pants cuffs below are essential curves, not specialized
peripheral trace constants. There are two concrete ways to advance the
preceding construction.
First, an explicit bounded coefficient functional cancels every
one-self-intersection primitive geodesic together with all its iterates,
simultaneously over all pants. Second, a product in the length-convolution
algebra gives a unit first atom for every weighted multicurve, including
those with disconnected support. Neither result proves a new counting
remainder. They also reveal why a fixed linear functional cannot remove
every nonsimple primitive curve.

\subsection{An explicit filter for the entire one-crossing sector}

Write $n(C)=\#C$, counting parallel copies, and $n(\varnothing)=0$.
This differs from the number $r(C)$ of distinct supports in
Section~\ref{subsec:countsupport}. We seek a functional
\[
 \Lambda(T_C)=w_{n(C)}
\]
which preserves simple primitive curves and annihilates the complete
iterate family of every primitive figure-eight.

A filter is a linear rule assigning a number to each basis
multicurve. It can make the contributions of a chosen geometric
family cancel when their trace polynomials are expanded.
For a figure-eight, $T_\gamma=xy-z$, so the simplest cancellation
requires $w_2-w_1=0$. Cancelling every power imposes more
conditions, which the following moment construction solves at once.

For $0\leq\theta\leq\pi$, put $v=2\cos\theta$ and let
$u_+,u_-$ be the two roots of
\[
 u^2-u=v.
\]
The unordered pair of roots is unambiguous, including where the
discriminant vanishes. Define
\begin{equation}
 w_n=\frac1\pi\int_0^\pi
       \bigl(u_+(\theta)^n+u_-(\theta)^n\bigr)\,\dd\theta,
 \qquad n\geq0.
 \label{eq:globalmoments}
\end{equation}
In particular $w_0=2$ and $w_1=1$. If $v\geq-1/4$, the roots are
real and lie in $[-1,2]$. If $v<-1/4$, they are complex conjugates
of modulus $\sqrt{-v}\leq\sqrt2$. Thus
\begin{equation}
 |w_n|\leq2^{n+1},\qquad
 |\Lambda(f)|\leq2\|f\|_2.
 \label{eq:globalbounded}
\end{equation}
Here $\|\cdot\|_2$ is the weighted coefficient norm from
Section~\ref{sec:coefficientdistributions}, not a norm obtained
by evaluating characters.

The weights have a closed generating function. The root sums
$S_n(v)=u_+^n+u_-^n$ satisfy
$S_0=2$, $S_1=1$, and $S_n=S_{n-1}+vS_{n-2}$, whence
\begin{equation}
 \boxed{\quad
 \sum_{n\geq0}w_n q^n
 =\frac{2-q}{\sqrt{(1-q)^2-4q^4}},
 \qquad |q|<\frac12.
 \quad}
 \label{eq:globalweightseries}
\end{equation}
The square root has value one at zero. To derive the formula,
average $(2-q)/(1-q-vq^2)$ over $v=2\cos\theta$ and use the
elementary integral of $(a-b\cos\theta)^{-1}$.
The first weights are
\[
 2,\ 1,\ 1,\ 1,\ 5,\ 11,\ 19,\ 29,\ 53,\ldots.
\]
They are nonnegative integers: the coefficients of $S_n(v)$ are
nonnegative integers and the averages of $v^{2j}$ and $v^{2j+1}$
are $\binom{2j}{j}$ and zero, respectively.

\begin{proposition}[A global one-crossing filter]
\label{prop:globalonecrossing}
For every simple primitive curve $\alpha$,
$\Lambda(T_\alpha)=1$. For every primitive geodesic $\gamma$
with exactly one transverse self-intersection and every $k\geq1$,
\begin{equation}
 \boxed{\Lambda(T_{\gamma^k})=0.}
 \label{eq:globalonecrossing}
\end{equation}
\end{proposition}

\begin{proof}
Use the pants neighborhood of a one-crossing geodesic and its
relation $T_\gamma=xy-z$, where $x,y,z$ are the unnegated cuff
traces~\cite[Theorem 1.1]{RivinFigures}. In a closed surface these
cuffs are essential; when two boundary curves have the same
geodesic representative, keep them distinct until their traces
are identified. Replacing every trace monomial of total
degree $n$ by $u^n$ sends $P_k(xy-z)$ to $P_k(u^2-u)$.
This remains true if two cuff images coincide: the identification
preserves total degree, including parallel multiplicities. Since
$u_\pm^2-u_\pm=2\cos\theta$,
\[
 \Lambda(T_{\gamma^k})
 =\frac2\pi\int_0^\pi P_k(2\cos\theta)\,\dd\theta
 =\frac4\pi\int_0^\pi\cos(k\theta)\,\dd\theta=0.
\]
The simple-curve assertion is $w_1=1$.
\end{proof}

The first two cases check the mechanism directly:
$\Lambda(xy-z)=w_2-w_1=1-1=0$, and
\[
 \Lambda((xy-z)^2-2)=w_4-2w_3+w_2-2w_0
                  =5-2+1-4=0.
\]
The last constant contributes $-2w_0=-4$ because
$\Lambda(1)=2$. By contrast a simple curve's double traversal
has value $\Lambda(x^2-2)=w_2-2w_0=-3$.
Thus the filter cancels the intended figure-eight family without
silently removing all iterates or producing a positive measure.

This is a single global functional, so no choice of a pants
decomposition and no sum of separately normalized pants projections
is required. It is not evaluation at a flat representation having
the same trace on every simple curve; no such assertion is used.

Equation~\eqref{eq:coefffullheatexpansion} and
\eqref{eq:globalbounded} justify the convergent aggregate
\begin{equation}
 \mathcal H_X^\Lambda(t)=\sum_C w_{n(C)}H_C^X(t),
 \qquad \Re t>0.
 \label{eq:globalfilteredheat}
\end{equation}
It is holomorphic in that half-plane. Since $\Lambda(1)=2$, the
constant identity contribution must also be multiplied by two if
one applies the functional before subtracting that contribution.
Absolute convergence of the
orbital weighted-norm estimate permits applying $\Lambda$ before
the orbital sum. Every one-crossing primitive and all its powers
therefore disappear. Simple iterates remain with the explicit
coefficients $\Lambda(P_k(T_\alpha))$; their first values are
$1,-3,-2,5,11,-6,\ldots$. Primitives with at least two
self-intersections also remain.

If the higher symmetric-power twisting spectra used for primitive
extraction in Section~\ref{sec:coefficientdistributions} are supplied,
let
\[
 \Pi_\rho^X=\sum_{\delta\in\mathcal P_{\rm un}}
                T_\delta(\rho)\delta_{\ell_\delta}.
\]
This is the primitive measure from that section divided by its
positive length variable. Then the filtered primitive measure satisfies
\begin{equation}
 \Lambda(\Pi_\rho^X)
 =\sum_{\alpha\ {\rm simple}}\delta_{\ell_\alpha}
      +\mathcal R_X,
 \qquad
 \mathcal R_X=
 \sum_{\substack{\delta\in\mathcal P_{\rm un}\\i(\delta,\delta)\geq2}}
      \Lambda(T_\delta)\delta_{\ell_\delta}.
 \label{eq:globalresidual}
\end{equation}
Here $i(\delta,\delta)$ is the minimal self-intersection number.
The remainder is signed. Its Gaussian transforms converge by the
same norm estimate; no smaller-order counting bound for it has
been established.

\subsection{Why a fixed linear filter cannot finish the job}

There is an exact obstruction already in one pair of pants.
Choose embedded geodesic pants with three distinct essential cuffs,
let $a,b$ generate its fundamental group, and put
\[
 \gamma_1=a^2ba^{-1}b^{-1},\qquad
 \gamma_2=a^3b^2,\qquad
 \gamma_3=a^2bab.
\]
Then, as functions on the $\mathrm{SL}_2$ character variety,
\begin{equation}
 \boxed{T_a=T_{\gamma_1}+T_{\gamma_2}-T_{\gamma_3}.}
 \label{eq:globalobstruction}
\end{equation}
For a direct verification set $x=T_a$, $y=T_b$, $z=T_{ab}$ and
$k=x^2+y^2+z^2-xyz-2$. Cayley--Hamilton gives
\begin{align*}
 T_{\gamma_1}&=x(k-1),\\
 T_{\gamma_2}&=(x^2-1)yz-xy^2-x^3+3x,\\
 T_{\gamma_3}&=xz^2-yz-x.
\end{align*}
The indicated combination is $x$.

All three words are primitive in the geodesic sense of not being
proper powers: their homology vectors are $(1,0)$, $(3,2)$, and
$(3,2)$. They are nonperipheral cyclically reduced words in the
pants group, so their geodesics are nonsimple. Embedded geodesic
pants are convex; these intersections persist in the ambient
surface. An ambient proper root would traverse the same geodesic
image inside the pants and give a root in the pants group, which
the homology vectors exclude. Here the word primitive does not assert membership in a
free basis of the pants group.

Consequently, any linear functional $L$ satisfying $L(T_a)=1$
must satisfy
\begin{equation}
 \max_{1\leq j\leq3}|L(T_{\gamma_j})|\geq\frac13.
 \label{eq:globalobstructionbound}
\end{equation}
In particular, no linear functional on this character algebra
preserves every primitive simple trace and annihilates every
primitive nonsimple trace. The assertion applies also to
cutoff-dependent functionals once the same functional is applied
to all four labels. It does not exclude functionals depending
on each atom's length, nonlinear operations, or richer
representation data. Our explicit filter has values $-5,5,-1$
on these three words, in agreement with the identity.

\subsection{Retaining lengths: an Euler product with a unit diagonal}

Convolution adds lengths: an atom at $a$ times one at $b$ becomes
an atom at $a+b$. For a single primitive curve the inverse factor
below is the formal geometric series
$\mathbf1+T_\alpha\delta_{\ell_\alpha}
 +T_\alpha^2\delta_{2\ell_\alpha}+\cdots$.
For disjoint $\alpha,\beta$, multiplying two such factors gives
$T_{\alpha+\beta}\delta_{\ell_\alpha+\ell_\beta}$ with coefficient
one. There is no division by two: a multiset is chosen once from
the product, rather than ordered and then averaged.

To preserve the information discarded by a fixed character
functional, work in the convolution algebra of locally finite
length measures supported on $[0,\infty)$, with
$\delta_u*\delta_v=\delta_{u+v}$.
Its coefficients lie in the character algebra and its unit is
$\mathbf1=T_\varnothing\delta_0$. All products are first defined
by finite length cutoffs, so each retained coefficient is a finite
character polynomial. Define
\begin{equation}
 \mathbf E_X(\rho)=
 \prod_{\delta\in\mathcal P_{\rm un}}
 \bigl(\mathbf1-T_\delta(\rho)\delta_{\ell_\delta}\bigr)^{*-1}.
 \label{eq:globalEuler}
\end{equation}
This is a product of scalar trace geometric series, not the
determinant Selberg product. The positive systole makes every
length cutoff finite. Equivalently,
\[
 \mathbf E_X(\rho)=
 \sum_{\mathfrak m}
   \left(\prod_\delta T_\delta(\rho)^{m_\delta}\right)
   \delta_{\sum_\delta m_\delta\ell_\delta}
 =\sum_C T_C(\rho)\nu_C^X,
\]
where $\mathfrak m$ runs over finite multisets of primitive
unoriented geodesics, including the empty multiset.

\begin{proposition}[A unit first atom for every multicurve]
\label{prop:globaldiagonal}
For every integral weighted multicurve $C$, including the empty one,
\begin{equation}
 \boxed{
 \nu_C^X=\delta_{\ell_X(C)}+\tau_C^X,\qquad
 \operatorname{supp}\tau_C^X\subset(\ell_X(C),\infty).
 }
 \label{eq:globaldiagonal}
\end{equation}
The first atom is isolated.
\end{proposition}

\begin{proof}
Resolve the crossings of the union of geodesics in an input
multiset. Every resulting multicurve has total length at most
the input length. If some primitive component is nonsimple, or
two distinct components intersect, resolving a genuine transverse
crossing and tightening gives strict shortening, as in the earlier
length-ordering argument. Repeated components can be separated by
small perturbations; a transverse crossing already present in a
component or between distinct geodesics retains a strict shortening.

Equality is therefore possible only when all primitive components
are simple and mutually disjoint, allowing repeated copies.
Their product is then exactly $T_C$, and the unique multiset giving
this $C$ has coefficient one. The geometric series in
\eqref{eq:globalEuler} supplies no factorial. Local finiteness
isolates the first length. The empty multiset gives the empty
multicurve's unit atom at zero.
\end{proof}

\subsection{Absolute convergence and spectral input}

The fixed-coefficient bracelet estimate extends to multicomponent
diagrams. With $t(u)=2\cosh(u/2)$ and the same $K_{C,X}$ as before,
\[
 \left|[T_C]\prod_\delta T_\delta^{m_\delta}\right|
 \leq K_{C,X}\prod_\delta t(\ell_\delta)^{m_\delta}.
\]
Indeed, sign coherence and positive Fuchsian evaluation apply to
the entire diagram, whose absolute evaluation is that product.
It follows that
\begin{equation}
 \int e^{-sL}\,\dd|\nu_C^X|(L)
 \leq K_{C,X}
 \prod_{\delta\in\mathcal P_{\rm un}}
 \bigl(1-t(\ell_\delta)e^{-s\ell_\delta}\bigr)^{-1}
 <\infty
 \label{eq:globalEulerbound}
\end{equation}
for real
\begin{equation}
 s>s_X^*=
 \max\left\{\frac32,
       \frac{\log(2\cosh(\varepsilon/2))}{\varepsilon}\right\},
 \qquad \varepsilon=\operatorname{sys}(X).
 \label{eq:globalEulerhalfplane}
\end{equation}
The second condition makes every geometric-series ratio less
than one: $\log t(u)/u$ decreases in $u$. The first makes the
sum of those ratios converge using the elementary $O_X(e^R)$
primitive-orbit bound. This also proves coefficient holomorphy
for $\Re s>s_X^*$ and convergence of every coefficient Gaussian
transform for $\Re t>0$. For the latter claim fix any real
$s_0>s_X^*$ in~\eqref{eq:globalEulerbound}; on a compact subset
of $\Re t>0$, the factor
$\exp(-L^2\Re(1/t)/4+s_0L)$ is uniformly bounded. This supplies
an integrable majorant against $e^{-s_0L}\,\dd|\nu_C^X|$.
No norm convergence of the whole character-valued product is
needed for this fixed-coefficient assertion.

Recovering this product spectrally requires its higher twisting
data to be kept straight. For a representation $V$, write
$\Pi_V=\sum_\delta\chi_V(\delta)\delta_{\ell_\delta}$ and let
$D_m\delta_L=\delta_{mL}$. The precise reconstruction formula is
\begin{equation}
 \boxed{
 \log_*\mathbf E_X(\rho)
 =\sum_{m\geq1}\frac1m D_m\Pi_{\rho^{\otimes m}}.
 }
 \label{eq:globalEulerinput}
\end{equation}
Here $\log_*$ is the usual logarithm series with multiplication
replaced by convolution; it is well defined at each positive length
cutoff. For a factor $1-X$, the identity
$-\log(1-X)=\sum_{m\ge1}X^m/m$ explains the $1/m$ and the
addition of $m$ copies of its source length.

The tensor-power character is $T_\delta^m$. It must not be
replaced by $P_m(T_\delta)=\operatorname{tr}\rho(\delta^m)$.
Tensor powers decompose into finite sums of $\mathrm{SU}(2)$
symmetric powers, and the earlier Adams inversion recovers their
primitive measures from the corresponding higher-twist spectra.
Thus all those twisting spectra suffice for
\eqref{eq:globalEulerinput}. A single rank-two primitive trace
measure does not directly provide the required power sums at
coincident lengths.

\subsection{What these constructions settle, and what remains}

The global filter gives an actual convergent aggregate with the
entire one-crossing sector removed. Equation~\eqref{eq:globalobstruction}
proves that eliminating all nonsimple primitives by a fixed linear
character functional is impossible. The length-valued product
avoids that obstruction: its first-atom extraction is nonlinear
and retains the source lengths.

In particular, the positive support partition function of
Section~\ref{subsec:countsupport} has the exact reconstruction
\[
 \mathcal P_X(s,z)
 =\sum_C z^{r(C)}
   \exp\!\left[-s\,\inf\operatorname{supp}\nu_C^X\right],
 \qquad \Re s>0.
\]
The unit first atoms make this identity unambiguous. It is a
reconstruction statement, not a remainder estimate. Quantitative
counting now requires either control of the residual
$\mathcal R_X$ in~\eqref{eq:globalresidual}, using cancellations
among higher-intersection curves, or uniform control of the
nonlinear diagonal extraction from $\mathbf E_X$. Neither follows
from fixed-coefficient convergence alone.

\paragraph{Finite computational companion.}
The maintained recipe \path{experiments/counting_mixing/run.py} checks
the trace recurrences through power $12$, four rational $SL_2$ word
fixtures, the global weights, finite support selectors, and a formal
length-valued Euler product through length $14$. Its result is
\path{artifacts/counting_mixing/result.json}; the input ranges and
independent review are recorded with the recipe and the checkpoint
\path{docs/execution/counting-mixing-review-2026-11-01/}.
These are exact finite algebraic checks. The formal lengths are not
asserted to come from one hyperbolic metric, and the checks do not
establish the infinite-product convergence, geometric first-atom
theorem, or a sharp counting remainder.

\clearpage

